\subsection{ARCH-136-MOORE-TOWER: a complete $p$-local one-disk tower}

\paragraph{Classical context and verification scope.}
This is a concrete cochain-graded verification of the one-generator
Cartan-resolution pattern; no historical novelty is claimed for that
resolution. For odd $p$, see Pedro Real, \emph{Homological perturbation
theory and associativity}, Homology, Homotopy and Applications
\textbf{2} (2000), 51--88, Theorem~5.28, p.~76, and
Theorem~5.32, p.~81. The former gives the corresponding reduced-bar
generators and $p$-Moore differential, and the latter promotes it
to a free Cartan resolution. If $r$ is divided-power arity and
$h$ the homological degree in that construction with its original
generator in degree two, the regrading $q=4r-h$ gives the degrees
$2p^i-1$ and $2p^i-2$ below. The proof here also covers $p=2$
and checks the exact first surviving class at every finite stage,
the weight bookkeeping, and stabilization of low form degrees.
These facts do not supply the separate target-compatible
coefficient-tail completion.

\paragraph{Statement.}
Retain $R=\mathbf Z_{(p)}$ and the initial disk $da=b$ of
ARCH-136-MOORE-PAIR.  Put $c_0=a$ and, for $r\geq0$, define the unit
\[
                 u_r=\frac1p\binom{p^{r+1}}{p^r}\in R^\times.
\]
For $r\geq1$ adjoin odd $c_r$ and even $t_r$ in form degrees
$2p^r-1$ and $2p^r-2$, both of weight $p^rm$, with
\[
 dc_r=b^{[p^r]},\qquad
 dt_r=pc_r-u_{r-1}^{-1}c_{r-1}
                            b^{[p^r-p^{r-1}]}.           \tag{136.MT1}
\]
Let $D_p^{(r)}$ be the algebra after stage $r$, where stage zero
is the initial disk.  Then
\[
 \begin{split}
 H^j(D_p^{(r)})&=0\quad(1\leq j\leq2p^{r+1}-1),\\
 H^{2p^{r+1}}(D_p^{(r)})&=R/pR,
 \end{split}                                             \tag{136.MT2}
\]
with the second group generated by $b^{[p^{r+1}]}$ in weight
$p^{r+1}m$.  The union $D_p^{(\infty)}$ has cohomology $R$ in
degree zero and zero in every positive degree.  Its low-degree
algebra stabilizes because new generator degrees tend to infinity.

\paragraph{Proof.}
The factorial valuation formula gives
$v_p\binom{p^{r+1}}{p^r}=1$, so $u_r$ is a unit.  The derivative
of the second term of (136.MT1) is
\[
 u_{r-1}^{-1}b^{[p^{r-1}]}b^{[p^r-p^{r-1}]}
      =u_{r-1}^{-1}\binom{p^r}{p^{r-1}}b^{[p^r]}
      =p b^{[p^r]}.
\]
Thus $d^2=0$, and all differentials preserve weight.

Use the variables $z_j=b^{[p^j]}$ of (136.MP4).  Modulo $p$, the
factorial digit formula writes
$b^{[p^{j+1}-p^j]}$ as a nonzero scalar multiple of $z_j^{p-1}$.
Consequently $D_p^{(r)}\otimes\mathbf F_p$ is an actual tensor
product of the factors
\[
 \Lambda(c_j)\otimes\mathbf F_p[z_j]/(z_j^p)
                    \otimes\Gamma(t_{j+1})\quad(0\leq j<r),
\]
with $dc_j=z_j$ and $dt_{j+1}$ a nonzero scalar multiple of
$c_jz_j^{p-1}$, followed by
$\Lambda(c_r)\otimes\mathbf F_p[z_r]/(z_r^p)$ and the remaining
zero-differential factors with indices $j>r$.
Each factor containing a $t_{j+1}$ contracts onto constants by
the explicit unit-pair contraction in ARCH-136-MOORE-PAIR.
The factor without its next $t$ has cohomology generated by
$1$ and $c_rz_r^{p-1}$.  Hence
\[
 H^*(D_p^{(r)}\otimes\mathbf F_p)
  =\Lambda(c_rz_r^{p-1})\otimes
                   \bigotimes_{j\geq r+1}
                              \mathbf F_p[z_j]/(z_j^p).   \tag{136.MT3}
\]
The first positive degrees are $2p^{r+1}-1$ and $2p^{r+1}$.

Rationally, make the triangular generator changes
\[
 c'_j=c_j-\frac1p u_{j-1}^{-1}c_{j-1}
                              b^{[p^j-p^{j-1}]}.
\]
They are invertible successively, and satisfy $dc'_j=0$ and
$dt_j=pc'_j$.  The algebra is the tensor product of the initial
rational odd unit disk and even unit disks, so it has only constant
cohomology.  At a finite stage every form-degree term is finite
free over $R$.  Apply the coefficient exact sequence (136.MP6) to
(MT3): as in the one-pair proof it gives all the vanishings in
(136.MT2).  The surviving closed class $b^{[p^{r+1}]}$ generates
the next integral group modulo $p$, and is killed by $p$ because
\[
 d\!\left(u_r^{-1}c_r b^{[p^{r+1}-p^r]}\right)
                          =p b^{[p^{r+1}]}.
\]
Finite generation and Nakayama's lemma give exactly $R/pR$.
For $r=0$ the same argument is the initial odd-disk calculation.

Filtered colimits of modules commute with cohomology.  Each fixed
positive degree is zero at every sufficiently large stage by
(136.MT2), so it is zero in the union.  The degree-zero constants
are unchanged.  Equally, one may observe that in each fixed finite
form range the underlying algebra itself becomes constant once
all later generators have larger degree.\hfill$\square$

\paragraph{The exact limitation.}
The tower repairs native cycles at weights $mp^r$.  Admissibility
of a target supplies a cutoff $\lambda_C(2p^r)$ for that degree,
with no required comparison between these two sequences.
Consequently the displayed tower cannot automatically retain an
arbitrary prescribed map from the initial disk to that target.
The next example shows that the issue persists for the guarded
free horizontal source used in this run; it is not removed by its
horizontal lifting property.

\subsection{ARCH-136-NATIVE-LIFT: guarded free sources need not lift native Moore primitives}

\paragraph{Statement.}
For every $m\geq2$ there is an admissible integral theory $C$ with
horizontal cutoffs $\mu_C(q)=q+1$, and a pure horizontal free source
$H\to C$ as in ARCH-136-DISK-SUPPORT with those cutoffs, having a
selected free odd unit-disk generator $x\in H_1^{m,1}$, such that
for every prime $p$ the image of
$\gamma_p(dx)$ is not a vertical boundary in $C_1^{mp,*}$.
In particular no forward extension of $H$ adjoining a primitive
of that native cycle can retain this prescribed map to $C$.
This refutes native-weight lifting, not coefficient-tail repair
or admissible common sources.

\paragraph{Proof.}
Let $U$ be the simplicial free abelian module corresponding under
Dold--Kan to the horizontal unit disk
\[
                \mathbf Z[2]\xrightarrow{1}\mathbf Z[1].
\]
It is simplicially contractible, $U_0=0$, and $U_1=\mathbf Z$.
Make two copies $V,W$ in form degrees one and two, both of weight
$m$, and form
\[
 E=\Lambda(V)\otimes\Gamma(W),\qquad dv=w,
 \qquad C=B\mathbin{\times}_{\mathbf Z}E.                 \tag{136.NL1}
\]
Both augmentations kill their positive-weight generators; the
coefficient map to $C$ is $(S\to B,\ S\to\mathbf Z\to E)$.
The fibre product and its differential, PD operations and
simplicial structure are strict.

As a bigraded module, $C=B\oplus I_E$.  In every positive weight,
each row of $E$ is a finite direct sum of functors
$\Lambda^i(U)\otimes\Gamma^j(U)$ of positive total arity.
Apply those functors to a linear simplicial contraction of $U$;
functorial simplicial homotopies give a contraction of each such
row.  Its zero map remains zero because the arity is positive.
Therefore $I_E$ contributes no horizontal homology at any weight.
The standard $B$ bound $w>q$ proves the claimed cutoffs for $C$.

In a fixed form degree $q$, a monomial in $E$ has weight
$m(i+j)$ with $i+2j=q$, hence weight at most $mq$.
Thus every extra vertical cohomology group in that degree has
weight at most $mq$, independently of the simplex dimension.
Also $E^0=\mathbf Z$ and has no extra degree-zero cohomology.
The standard vertical comparison for $B$ consequently proves
admissibility of $C$.  Notice that its extra fixed-weight vertical
defects are allowed by the actual definition.

In the free horizontal construction choose its zero-boundary cell
in form degree one, weight $m$, and horizontal dimension one.
It is allowed because $m\geq\mu_C(1)=2$.  Its normalized generator
can be sent to the lower generator of $U$ in horizontal degree one:
that generator has zero horizontal boundary.  Send its vertical
partner to the corresponding generator of $W$.
The free property gives a strict target map at this stage.  Complete
the construction of ARCH-136-DISK-SUPPORT, choosing all further
horizontal primitives in $C$; its acyclic rows permit every step
without changing this initial choice.  Call the resulting source
$H$ and its selected generator $x$.

In simplex one the extra algebra is exactly
$E_1=\Lambda(v)\otimes\Gamma(w)$ with $dv=w$.
The class of $w^{[p]}$ is the nonzero element of
$H^{2p}(E_1^{mp,*})=\mathbf Z/p$ from (136.DS4).
It survives in $C_1$ because the extra augmentation ideal is a
direct summand of the underlying vertical complex.  It is precisely
the image of $\gamma_p(dx)$.  A strict map carries a boundary to
a boundary, so adjoining its native primitive cannot extend the
chosen map.  This proves the assertion.\hfill$\square$

\paragraph{Why a tail repair is still distinct.}
In this example every positive-weight coefficient acts as zero on
$I_E$.  Thus the positive coefficient multiples of the obstructing
native class already have zero image in that summand.  The example
therefore gives no obstruction to repairing sufficiently high
coefficient multiples, exactly as required by admissibility.
Doing so coherently introduces the relations of the coefficient
PD algebra and the repeated-presentation issue of
ARCH-136-MOORE-DUPLICATE; neither is settled by the one-disk tower.
